\documentclass[12pt,oneside,a4paper,chapter=TITLE]{abntex2}

\usepackage[T1]{fontenc}
\usepackage[utf8]{inputenc}
\usepackage[brazil]{babel}
\usepackage{indentfirst}
\usepackage{graphicx}
\usepackage{lmodern}
\usepackage{booktabs}
\usepackage{float}
\usepackage{url}
\usepackage{hyperref}

\setlength{\parindent}{1.25cm}
\setlength{\parskip}{0.2cm}

\titulo{UniSync}

\autor{
Roberto Matheus Nunes dos Santos \\
Kamilly Stephany Saraiva Santos \\
João Victor Soares de Lima \\
Samuel Alexandre Correia da Silva \\
Reinoldes Clarindo de Oliveira Junior
}

\instituicao{Centro Universitário Maurício de Nassau - UNINASSAU}
\local{Maceió - AL}
\data{2026}
\tipotrabalho{Relatório Técnico}

\begin{document}

\begin{figure}[!htb]
    \centering
    \includegraphics[width=0.35\textwidth]{logo_unina.png}
\end{figure}

\imprimircapa

\chapter{Introdução}

O \textbf{UniSync} é a proposta de uma plataforma voltada à centralização
de oportunidades acadêmicas, como palestras, eventos, estágios e
projetos. A ideia surge a partir da necessidade de facilitar o acesso
dos estudantes a informações que podem estar distribuídas em diferentes
canais e ambientes.

A proposta consiste na organização dessas informações em uma solução
integrada, permitindo que o estudante encontre oportunidades de
interesse de forma mais organizada e tenha acesso a informações
relacionadas a prazos e notificações.

Para atender a essa proposta, o UniSync deverá ser estruturado com uma
arquitetura baseada em microsserviços e Micro-Frontends, possibilitando
a separação dos diferentes componentes da solução.

O projeto está organizado em componentes para ambiente Web, serviços de
backend, aplicativo mobile e sistemas de persistência de dados. Cada
parte deverá possuir responsabilidades específicas, permitindo uma
organização modular da solução.

\chapter{Contextualização do Problema}

Informações relacionadas a oportunidades acadêmicas podem estar
disponíveis em diferentes meios de comunicação, dificultando sua
localização e acompanhamento pelos estudantes.

Palestras, eventos, estágios e projetos podem possuir diferentes
períodos de inscrição, datas e informações específicas. Quando esses
dados estão distribuídos em diferentes canais, o estudante pode ter
dificuldade para acompanhar as oportunidades disponíveis e seus
respectivos prazos.

Nesse contexto, o UniSync é proposto como uma plataforma capaz de
centralizar essas informações em um único ambiente, facilitando a
consulta e o acompanhamento das oportunidades acadêmicas.

\chapter{Justificativa}

A proposta do UniSync justifica-se pela possibilidade de oferecer aos
estudantes um ambiente centralizado para consulta de oportunidades
acadêmicas.

A utilização de uma arquitetura baseada em microsserviços e
Micro-Frontends também permite propor uma solução modular, na qual
diferentes partes da plataforma possuem responsabilidades bem
definidas.

Além disso, a existência de uma aplicação mobile poderá facilitar o
acesso às informações e ao acompanhamento de prazos por meio de
dispositivos Android.

Dessa forma, o projeto busca propor uma solução tecnológica organizada
para facilitar o acesso às oportunidades acadêmicas e melhorar a
disponibilidade dessas informações para os estudantes.

\chapter{Objetivos}

\section{Objetivo Geral}

Propor uma plataforma denominada \textbf{UniSync} para centralização de
oportunidades acadêmicas, utilizando uma arquitetura baseada em
microsserviços e Micro-Frontends.

\section{Objetivos Específicos}

\begin{itemize}
    \item Centralizar informações sobre palestras, eventos, estágios e
    projetos;

    \item Disponibilizar um feed para apresentação das oportunidades;

    \item Permitir a busca e filtragem das oportunidades;

    \item Possibilitar o acompanhamento de prazos relacionados às
    oportunidades;

    \item Prever a utilização de alertas em tempo real no aplicativo
    mobile;

    \item Organizar a solução por meio de microsserviços e
    Micro-Frontends;

    \item Propor uma aplicação mobile nativa para dispositivos Android.
\end{itemize}

\chapter{Visão Geral da Solução}

O UniSync deverá ser composto por diferentes componentes responsáveis
por partes específicas da solução. A arquitetura proposta contempla
uma aplicação Web, serviços de backend, um aplicativo mobile e
diferentes mecanismos de persistência de dados.

\section{Frontend Web}

O ambiente Web deverá utilizar uma arquitetura baseada em
Micro-Frontends. A proposta contempla um Shell responsável pela
estrutura principal da aplicação e um Micro-Frontend específico para
o feed de oportunidades.

O Shell será desenvolvido utilizando Python e Django, sendo responsável
pela estrutura do ambiente, roteamento principal e identidade visual
da aplicação.

O Micro-Frontend do feed deverá utilizar Vue.js, proporcionando uma
interface reativa para apresentação das oportunidades acadêmicas.

\section{Backend}

O backend deverá ser dividido em diferentes serviços, seguindo a
proposta de arquitetura de microsserviços.

O \textbf{Núcleo Acadêmico}, desenvolvido em \textbf{C\#} com .NET,
deverá ser responsável pelo gerenciamento de usuários, permissões,
autenticação e regras estruturadas relacionadas ao sistema.

O \textbf{Motor de Oportunidades}, desenvolvido em Python com FastAPI,
deverá fornecer uma API voltada às operações relacionadas às
oportunidades, incluindo mecanismos para consultas, buscas e
filtragens.

\section{Aplicativo Mobile UniSync}

O projeto também deverá contemplar um aplicativo mobile denominado
\textbf{UniSync}, desenvolvido em Kotlin utilizando Jetpack Compose.

A aplicação deverá ser direcionada a dispositivos Android e deverá
permitir o acompanhamento das oportunidades acadêmicas, incluindo
informações relacionadas a prazos e a possibilidade de utilização de
alertas em tempo real.

\chapter{Arquitetura Proposta}

\section{Organização dos Componentes}

A arquitetura do UniSync deverá seguir uma abordagem modular, separando
as responsabilidades entre os componentes Web, mobile e backend.

A utilização de microsserviços deverá permitir a separação das
responsabilidades do backend, enquanto os Micro-Frontends deverão
possibilitar a organização independente de partes da interface Web.

\section{Shell e Micro-Frontend}

O Shell deverá funcionar como a estrutura principal do ambiente Web,
sendo responsável pelo roteamento e pela organização da aplicação.

O Micro-Frontend desenvolvido em Vue.js deverá concentrar a
funcionalidade relacionada ao feed de oportunidades.

Essa separação deverá permitir que diferentes partes da interface
possuam responsabilidades específicas dentro da arquitetura proposta.

\section{Microsserviços}

O backend deverá ser organizado em serviços independentes.

O Núcleo Acadêmico deverá concentrar funcionalidades relacionadas aos
usuários, autenticação, permissões e regras estruturadas.

O Motor de Oportunidades deverá concentrar funcionalidades relacionadas
à consulta, busca e filtragem das oportunidades acadêmicas.

\section{Comunicação entre os Componentes}

A comunicação entre os componentes deverá ocorrer por meio de APIs,
permitindo que as diferentes partes da solução possam trocar
informações de maneira organizada.

No ambiente Web, o Shell deverá organizar a aplicação e realizar a
integração com o Micro-Frontend responsável pelo feed. O
Micro-Frontend deverá apresentar as oportunidades ao usuário e utilizar
os serviços disponibilizados pelo backend.

O aplicativo mobile UniSync também deverá se comunicar com os serviços
de backend, permitindo o acesso às informações relacionadas às
oportunidades, prazos e alertas previstos.

No backend, o Núcleo Acadêmico e o Motor de Oportunidades deverão
possuir responsabilidades distintas. O Núcleo Acadêmico deverá
concentrar informações relacionadas a usuários, autenticação,
permissões e regras estruturadas, enquanto o Motor de Oportunidades
deverá concentrar operações relacionadas à consulta, busca e filtragem
das oportunidades.

Os serviços de backend deverão acessar as estruturas de persistência
conforme o tipo de informação armazenada. Os dados estruturados,
incluindo usuários e turmas, deverão utilizar um banco relacional,
enquanto informações relacionadas a eventos, vagas e oportunidades
poderão utilizar o MongoDB.

De forma conceitual, a comunicação entre os componentes poderá ser
representada pelos seguintes fluxos:

\begin{center}
\textbf{Shell Web}
$\leftrightarrow$
\textbf{Micro-Frontend Feed}
$\leftrightarrow$
\textbf{APIs do Backend}
\end{center}

\begin{center}
\textbf{UniSync Mobile}
$\leftrightarrow$
\textbf{APIs do Backend}
\end{center}

\begin{center}
\textbf{Núcleo Acadêmico}
$\leftrightarrow$
\textbf{Motor de Oportunidades}
$\leftrightarrow$
\textbf{Camada de Persistência}
\end{center}

Esses fluxos representam a comunicação conceitual entre os
componentes e não correspondem a uma implementação concluída, uma vez
que o presente trabalho apresenta a ideação da solução.

\chapter{Tecnologias Propostas}

A arquitetura do UniSync é composta por diferentes camadas que utilizam
tecnologias específicas de acordo com suas responsabilidades. A tabela
a seguir apresenta a organização tecnológica proposta para cada
componente.

\begin{table}[H]
\centering
\caption{Tecnologias propostas para o UniSync}
\begin{tabular}{lll}
\toprule
\textbf{Camada} & \textbf{Componente} & \textbf{Tecnologia} \\
\midrule
Frontend Web & Shell & Python + Django \\
Frontend Web & Micro-Frontend Feed & Vue.js \\
Backend & Núcleo Acadêmico & C\# + .NET \\
Backend & Motor de Oportunidades & Python + FastAPI \\
Mobile & UniSync & Kotlin + Jetpack Compose \\
Persistência & Banco relacional & SQL Server ou PostgreSQL \\
Persistência & Banco não relacional & MongoDB \\
\bottomrule
\end{tabular}
\end{table}

\section{Python e Django}

Python será utilizado no desenvolvimento do Shell Web por meio do
framework Django. Essa camada deverá ser responsável pela estrutura
principal do ambiente Web.

\section{Vue.js}

Vue.js deverá ser utilizado no desenvolvimento do Micro-Frontend
responsável pelo feed de oportunidades.

\section{C\# e .NET}

O Núcleo Acadêmico deverá utilizar \textbf{C\#} e .NET para implementar
as funcionalidades relacionadas aos usuários, autenticação, permissões
e regras estruturadas.

\section{Python e FastAPI}

Python com FastAPI deverá ser utilizado no Motor de Oportunidades,
fornecendo uma API voltada às operações de consulta, busca e filtragem.

\section{Kotlin e Jetpack Compose}

O aplicativo mobile UniSync deverá ser desenvolvido utilizando Kotlin
e Jetpack Compose, com foco em dispositivos Android.

\section{Bancos de Dados}

A proposta contempla a utilização de um banco de dados relacional,
como SQL Server ou PostgreSQL, para armazenamento de informações
estruturadas, incluindo usuários e turmas.

O MongoDB deverá ser utilizado para informações relacionadas a eventos,
vagas e outras oportunidades que possam apresentar estruturas mais
flexíveis.

A definição final do modelo de dados e dos relacionamentos deverá
ocorrer durante as etapas posteriores de desenvolvimento.

\chapter{Funcionalidades Idealizadas}

\section{Centralização de Oportunidades}

O UniSync deverá centralizar informações relacionadas a palestras,
eventos, estágios e projetos em um único ambiente.

A proposta é permitir que o estudante tenha acesso a diferentes tipos
de oportunidades sem a necessidade de consultar diversos ambientes
separadamente.

\section{Feed de Oportunidades}

A aplicação Web deverá possuir um feed destinado à apresentação das
oportunidades disponíveis.

O feed deverá utilizar Vue.js por meio de um Micro-Frontend, permitindo
uma interface voltada à apresentação dinâmica das informações.

\section{Busca e Filtragem}

O sistema deverá contemplar mecanismos de busca e filtragem das
oportunidades.

A finalidade prevista é permitir que o estudante encontre informações
de acordo com os critérios disponíveis para cada oportunidade.

\section{Acompanhamento de Prazos e Alertas}

O aplicativo \textbf{UniSync} deverá contemplar o acompanhamento de
prazos relacionados às oportunidades. A finalidade prevista é permitir
que o estudante tenha acesso a informações temporais importantes
diretamente pelo aplicativo.

Além do acompanhamento de prazos, também está prevista a possibilidade
de utilização de alertas em tempo real no UniSync. Esses alertas poderão
auxiliar na apresentação de informações relacionadas às oportunidades e
aos seus respectivos prazos.

A implementação detalhada desse mecanismo deverá ser definida em uma
etapa posterior do desenvolvimento, uma vez que o presente trabalho
representa apenas a ideação do projeto.

\chapter{Dados e Persistência}

A persistência dos dados deverá ser organizada de acordo com as
características das informações utilizadas pela plataforma.

Para dados estruturados, como usuários e turmas, está prevista a
utilização de um banco de dados relacional, podendo ser adotado SQL
Server ou PostgreSQL.

Para informações relacionadas às oportunidades, eventos e vagas, está
prevista a utilização do MongoDB, permitindo maior flexibilidade na
estrutura desses dados.

A definição final do modelo de dados e dos relacionamentos deverá
ocorrer durante as etapas posteriores de desenvolvimento.

\chapter{Estrutura Inicial do Projeto}

A organização inicial do repositório deverá seguir uma estrutura
separada por componentes, facilitando a organização do desenvolvimento.

Entre os principais diretórios previstos estão:

\begin{itemize}
    \item \texttt{.github/workflows/}: arquivos relacionados aos
    fluxos de trabalho automatizados do repositório;

    \item \texttt{backend-python/}: componente relacionado aos serviços
    desenvolvidos em Python;

    \item \texttt{frontend-shell/}: componente responsável pelo Shell
    da aplicação Web;

    \item \texttt{frontend-web-aluno/}: componente relacionado à
    interface Web destinada ao aluno;

    \item \texttt{mobile-kotlin/}: componente destinado ao aplicativo
    mobile UniSync;

    \item \texttt{docs/}: documentação do projeto;

    \item \texttt{README.md}: documentação inicial e informações
    gerais do repositório.
\end{itemize}

Essa organização deverá facilitar a separação dos diferentes
componentes previstos na arquitetura do projeto.

\chapter{Organização da Equipe}

A divisão inicial das responsabilidades da equipe está apresentada a
seguir:

\begin{itemize}
    \item \textbf{Matheus}: Backend e Banco de Dados;

    \item \textbf{Kamilly}: Frontend;

    \item \textbf{João}: Backend e Banco de Dados;

    \item \textbf{Samuel Alexandre Correia da Silva}: Backend e
    Documentação em LaTeX;

    \item \textbf{Reinoldes}: Frontend.
\end{itemize}

A distribuição das atividades poderá ser ajustada de acordo com as
necessidades identificadas nas etapas posteriores do projeto.

\chapter{Etapas Futuras do Projeto}

Como o presente trabalho representa a ideação do UniSync, as etapas
seguintes deverão contemplar o detalhamento e desenvolvimento da
solução proposta.

Entre as atividades futuras poderão estar:

\begin{itemize}
    \item detalhamento dos requisitos funcionais e não funcionais;

    \item definição dos contratos das APIs;

    \item modelagem dos bancos de dados;

    \item desenvolvimento dos microsserviços;

    \item desenvolvimento dos componentes Web;

    \item desenvolvimento do aplicativo mobile;

    \item integração entre os diferentes componentes;

    \item realização de testes;

    \item avaliação da solução proposta.
\end{itemize}

Essas etapas deverão ser realizadas de forma progressiva, permitindo
que a arquitetura e as funcionalidades inicialmente propostas sejam
detalhadas e desenvolvidas.

\chapter{Considerações Finais}

O UniSync apresenta uma proposta de plataforma para centralização de
oportunidades acadêmicas, reunindo informações relacionadas a
palestras, eventos, estágios e projetos em um ambiente organizado.

A arquitetura baseada em microsserviços e Micro-Frontends foi proposta
com o objetivo de separar responsabilidades entre os diferentes
componentes da solução, contemplando ambientes Web, backend e mobile.

O aplicativo mobile UniSync deverá complementar a plataforma,
possibilitando o acompanhamento de oportunidades e de seus respectivos
prazos, além da possibilidade de utilização de alertas em tempo real.

Por se tratar de uma proposta de projeto, os detalhes de implementação,
integração e validação deverão ser definidos nas etapas futuras de
desenvolvimento.

\chapter{Referências}

\begin{itemize}
    \item DJANGO SOFTWARE FOUNDATION. \textit{Django documentation}.
    Disponível em: \url{https://docs.djangoproject.com/}. Acesso em:
    2026.

    \item VUE.JS. \textit{Vue.js documentation}. Disponível em:
    \url{https://vuejs.org/}. Acesso em: 2026.

    \item MICROSOFT. \textit{.NET documentation}. Disponível em:
    \url{https://learn.microsoft.com/dotnet/}. Acesso em: 2026.

    \item FASTAPI. \textit{FastAPI documentation}. Disponível em:
    \url{https://fastapi.tiangolo.com/}. Acesso em: 2026.

    \item ANDROID DEVELOPERS. \textit{Jetpack Compose}. Disponível em:
    \url{https://developer.android.com/compose}. Acesso em: 2026.

    \item MONGODB. \textit{MongoDB documentation}. Disponível em:
    \url{https://www.mongodb.com/docs/}. Acesso em: 2026.
\end{itemize}

\end{document}